% !TEX ROOT = ../Thesis.tex 
%-
%-> 中文摘要
%-
\chapter{摘\quad 要}\chaptermark{摘\quad 要}
\linespread{1.5}
\zihao{-4}

在大语言模型与检索增强生成（Retrieval-Augmented Generation, RAG）技术迅猛发展的背景下，多轮对话检索已成为对话式人工智能系统落地应用的关键环节。然而，真实多轮对话中普遍存在的指代、省略、话题漂移等语言现象，使得当前轮次的查询往往语义残缺，无法被现有检索器直接理解和召回。查询重写（Query Rewriting）作为在检索之前将上下文依赖的用户问句改写为语义自包含独立查询的关键预处理模块，直接决定了下游检索与生成的上限。然而，现有查询重写方法仍面临三大挑战：其一，对长历史、跨话题对话中的噪声轮次缺乏显式建模，导致重写时引入大量无关信息；其二，指代消解与省略补全被混合为端到端序列生成任务，训练目标不明确且错误难以定位；其三，重写目标局限于"像人写的参考重写"，与下游检索器的真实偏好存在系统性偏差。针对上述痛点，本文系统研究了面向多轮对话检索的查询重写方法，提出从"历史剪枝—渐进式重写—检索反馈优化"三阶段构建查询重写框架，并在多个公开基准上进行了深入的实证分析。本文的主要研究工作与创新成果如下：

首先，针对多轮对话中历史轮次的话题漂移问题，提出了一种话题漂移感知的对话历史剪枝方法。不同于将完整历史一视同仁地拼接输入的传统做法，本方法在轮次级引入三分类相关性判别器（强相关 / 弱相关 / 无关），基于 ModernBERT-base 构建轻量级分类头，并利用 TopiOCQA 的话题段落标注与 QReCC 的人工重写作为弱监督训练信号。在 94\,422 条轮次-查询对上训练后，分类器在留出验证集上取得 80.1\% 的准确率与 0.797 的宏 F1，能够稳定地压缩历史、降低输入噪声。

其次，针对指代消解与省略补全耦合导致错误难以分离的问题，提出了一种两阶段渐进式查询重写框架（Progressive Rewriter, PRW）。该框架将重写解耦为"指代消解"与"省略补全"两个串行子任务：第一阶段仅将指代词替换为对应实体，保持表层结构不变；第二阶段在此基础上补全缺失的主语、谓语与修饰成分，生成自包含查询。基于 Qwen2.5-3B-Instruct 与 LoRA（$r{=}16$）的两阶段 SFT 实验表明，PRW 在 TopiOCQA、Doc2Dial 两个话题切换密集、上下文依赖强的基准上，将 Recall@5 相对 No-Rewrite 分别提升 39 和 66 个百分点以上（BGE-large 检索器下）。

最后，针对重写目标与下游检索器偏好失配的问题，提出了检索反馈驱动的强化微调方法（Retrieval-Feedback Rewriter, RFR）。本文从 TopiOCQA 训练集的 10\,000 个带 gold passage 的轮次采样 80\,000 条候选重写，以 BGE-large 计算"候选 vs gold passage"余弦相似度作为反馈，构造 4\,106 条高质量偏好对，再用 Direct Preference Optimization (DPO) 在 PRW 基础上继续 LoRA 训练。在 129 步训练内，reward accuracies 从 45\% 单调上升至 85\%，训练全程稳定。在 QReCC、TopiOCQA、Doc2Dial、QuAC 四个公开基准 $\times$ BGE-large / GTE-ModernBERT / Dragon-Multiturn 三种异构检索器共 12 组设定中，RFR 相对 PRW 取得平均 $+1.10$ R@5 / $+1.31$ R@20 的一致增益，且该增益在训练从未见过的 GTE-ModernBERT 与 Dragon-Multiturn 上同样成立。

综上所述，本文围绕多轮对话检索中的查询重写问题，构建了一套"噪声过滤—结构化重写—下游反馈"相互衔接的技术方案，深入研究了对话历史建模、渐进式生成与基于检索反馈的偏好优化等核心问题，为工业级对话式 RAG 系统提供了一种轻量、高效且鲁棒的查询重写范式。

\keywords{多轮对话检索；查询重写；检索增强生成；直接偏好优化；大语言模型}

%-
%-> 英文摘要
%-
\chapter{Abstract}\chaptermark{Abstract}

Against the backdrop of the rapid progress of large language models (LLMs) and retrieval-augmented generation (RAG), multi-turn conversational retrieval has become a critical building block for deploying conversational AI systems. In realistic multi-turn dialogues, however, linguistic phenomena such as coreference, ellipsis, and topic drift are ubiquitous, rendering the current-turn query semantically incomplete and hard for existing retrievers to understand. Query rewriting (QR), which transforms a context-dependent utterance into a self-contained standalone query before retrieval, directly determines the upper bound of downstream retrieval and generation. Nevertheless, existing QR approaches still suffer from three key limitations: (i) they lack explicit modeling of noisy turns in long, topic-switching histories, so that irrelevant information is easily leaked into the rewrite; (ii) coreference resolution and ellipsis completion are conflated into a single end-to-end sequence generation task with ill-defined training signals; and (iii) the training objective is tied to "human-written references" and is systematically misaligned with the true preferences of downstream retrievers. To address these issues, this thesis systematically investigates query rewriting methods for multi-turn conversational retrieval, and proposes a three-stage framework consisting of \emph{history pruning}, \emph{progressive rewriting}, and \emph{retrieval-feedback optimization}. The main contributions are as follows.

First, to tackle topic drift in long conversational histories, we propose a \textbf{topic-drift-aware history pruning} module. The method introduces a turn-level three-way relevance classifier (strongly-, weakly-, and un-related) built on top of ModernBERT-base, trained with weak supervision derived from the topic-section labels of TopiOCQA and the human-written rewrites of QReCC. Trained on 94{,}422 turn--query pairs, the classifier achieves 80.1\% accuracy and macro-F1 = 0.797 on a held-out validation set, reliably compressing history and reducing input noise.

Second, to resolve the coupling between coreference and ellipsis phenomena, we propose a \textbf{Progressive Rewriter (PRW)} framework that decomposes rewriting into two sequential subtasks: a coreference-only stage that replaces anaphoric mentions while preserving surface structure, and an ellipsis-completion stage that restores missing subjects, predicates, and modifiers. With Qwen2.5-3B-Instruct fine-tuned via LoRA ($r{=}16$), PRW raises Recall@5 by more than 39 and 66 percentage points over No-Rewrite on TopiOCQA and Doc2Dial respectively (BGE-large retriever), i.e., the two benchmarks where topic switches and context dependencies are most severe.

Third, to bridge the gap between human-style rewrites and retriever-friendly rewrites, we design a \textbf{Retrieval-Feedback Rewriter (RFR)}. We sample 80{,}000 candidate rewrites from 10{,}000 TopiOCQA training turns with gold passages, score each candidate by its BGE-large cosine similarity to the gold passage, and assemble 4{,}106 high-quality preference pairs. DPO is then applied on top of PRW's LoRA adapter. Within only 129 optimization steps, reward accuracies rise monotonically from 45\% to 85\%. On 12 settings (four public benchmarks $\times$ three heterogeneous retrievers), RFR delivers a consistent average gain of $+1.10$ Recall@5 / $+1.31$ Recall@20 over PRW, and the gain transfers to unseen retrievers (GTE-ModernBERT and Dragon-Multiturn).

In summary, this thesis presents a unified technical route that couples noise filtering, structured rewriting, and downstream feedback, and provides a lightweight, efficient and robust query-rewriting paradigm for industrial-scale conversational RAG systems.

\englishkeywords{multi-turn conversational retrieval; query rewriting; retrieval-augmented generation; direct preference optimization; large language models}
